\section{An unresolved model with complete analytic control}\label{sec:variation}
Differentiating a parameter-dependent approximation to $\beta(2)$ can
produce one for $\beta(4)$: differentiation raises the denominator
power in the defining series. The difficulty is to cancel the remaining
$\beta(2)$ term and control the sign of the new error. A moment matrix
does both here. Its block determinant isolates a rational approximation,
and a variational formula shows that the differentiated error is positive.
This section therefore reaches the final arithmetic question rather
than proving another failure theorem: how quickly do the reduced
denominators of these rational approximations grow?

The ordinary-moment first variation \record{P07} is retained because
its analytic conclusions are stronger than a negative summary of the
project might suggest. It has all-index positivity and a sharp scalar
error rate. The remaining issue is primitive arithmetic.

The elementary Laplace identity
\[
 \int_0^\infty e^{-at}\frac{\sinh(\sqrt u\,t)}{\sqrt u}\dd t
       =\frac1{a^2-u}\qquad(a\geq1,\ |u|<1)
\]
suggests a parameter whose derivative raises the beta denominator
from the second to the fourth power. For $|u|<1$ let
\[
 F(u)=\sum_{k\geq0}\frac{(-1)^k}{(2k+1)^2-u},\qquad
 q_j(u)=\frac1{(2j+1)^2-u},\qquad
 w_u(x)=\frac{\sinh(\sqrt u(-\log x))}{\sqrt u(1+x^2)}.
\]
Use the analytic value of the quotient at $u=0$. The density is positive
for $0\leq u<1$; it need not be positive for negative $u$.
The moments $m_j(u)=\int_0^1x^{2j}w_u(x)\dd x$ obey
$m_j+m_{j+1}=q_j$, with $m_0=F$, $F(0)=\beta(2)$ and $F'(0)=\theta$.
These identities follow by expansion and Laplace integration, locally
uniformly in $u$.

Define $H_n(u)=\det[m_{i+j}(u)]_{0\leq i,j<n}$.
The identity $m_j+m_{j+1}=q_j$ suggests using multiples of $1+t$
as basis elements: their inner products with monomials are rational
functions of $u$. For $d=n-1$ set $c_j=q_j$ and
$B_{ij}=q_{i+j}+q_{i+j+1}$,
$0\leq i,j<d$. The same moment identities give
$B_{ij}=\int_0^1x^{2(i+j)}(1+x^2)^2w_u(x)\dd x$.
Thus $B(0)$ is positive definite and $B^{-1}$ is analytic near zero.
Put $a_n=\det B$ and $r_n=c^TB^{-1}c$; for $n=1$
use $a_1=1,r_1=0$. In the basis
$1,(1+t),(1+t)t,\ldots$, $t=x^2$, the moment matrix is
\[
 \begin{pmatrix}F&c^T\\c&B\end{pmatrix},\qquad
 H_n=a_n(F-r_n).
\]
The change of basis has determinant one; the formula for $H_n$
is the Schur-complement identity for the displayed block matrix.
The functions $a_n(u)$ and $r_n(u)$ are rational functions with rational
coefficients, regular at zero; their values and derivatives there are
rational. The following antisymmetric combination of the determinant
and its derivative eliminates $\beta(2)$ exactly:
\begin{equation}\label{eq:variation}
 W_n=a_n(0)H_n'(0)-a_n'(0)H_n(0)
       =a_n(0)^2(\theta-r_n'(0)).
\end{equation}
This is a formal identity, with no independence assumption about periods.
For example, at $n=2$ one obtains
$r_2(u)=q_0(u)^2/(q_0(u)+q_1(u))$ and $r_2'(0)=49/50$.
The resulting primitive form is $50\theta-49$.

\begin{proposition}[Positivity and the scalar rate; P07]\label{prop:variation}
Let $E_n=\theta-r_n'(0)$ and $\rho=3+\sqrt8$. Then $E_n>0$ for
every $n$, and
\[
 \lim_{n\to\infty}\frac{\log E_n}{n}=-2\log\rho.
\]
\end{proposition}
\begin{proof}
\emph{Positivity of the differentiated error.}
For $0\leq u<1$, completing the square in the block matrix gives
the following minimum over real polynomials:
\[
 F(u)-r_n(u)=
 \min_{\substack{\deg P<n\\P(-1)=1}}
                   \int_0^1P(x^2)^2w_u(x)\dd x.
\]
Every polynomial of degree less than $n$ with $P(-1)=1$ has the
form $1+(1+t)Q(t)$, $\deg Q<n-1$. Minimizing this quadratic form
in the coefficients of $Q$ gives the same $B^{-1}c$ as the block
determinant. The unique minimizing polynomial is
$P_{n,u}(t)=1-(1+t)\sum_{j<d}(B^{-1}c)_jt^j$ and is analytic
near zero. Differentiate from $u\geq0$ at zero. The derivative of
the minimizing polynomial satisfies $\dot P(-1)=0$, and stationarity
gives $\int P\dot P w_0\dd x=0$. Thus only the derivative of the
weight contributes. Expansion of $\sinh(\sqrt u(-\log x))/\sqrt u$
shows that this derivative is $(-\log x)^3/6$ divided by $1+x^2$.
Thus
\begin{equation}\label{eq:variation-positive}
 E_n=\int_0^1P_{n,0}(x^2)^2
                         \frac{(-\log x)^3}{6(1+x^2)}\dd x>0.
\end{equation}

\emph{The exponential rate.} The same polynomial minimizes the
weight-two norm, but its weight-four norm is the error we need.
A lower bound follows from the weight-four minimum. For an upper
bound we must control its supremum norm from its weight-two norm;
the ratio of the two densities is unbounded near zero, so a direct
pointwise comparison would not suffice.

Let $\lambda_{n,s}$ be the analogous constrained minimum for the
weight $(-\log x)^{s-1}/[(s-1)!(1+x^2)]$, $s=2,4$.
Then $E_n\geq\lambda_{n,4}$, while $P_{n,0}$ minimizes the
weight-two norm. In the variable $t=x^2$ the densities are
\[
 \eta_2(t)=\frac{-\log t}{4\sqrt t(1+t)},\qquad
 \eta_4(t)=\frac{(-\log t)^3}{96\sqrt t(1+t)}.
\]
The test polynomial $T_{n-1}(2t-1)/T_{n-1}(-3)$ gives
$\limsup n^{-1}\log\lambda_{n,s}\leq-2\log\rho$.

For the opposite bound, restrict to a fixed interior interval
$I_\epsilon=[\epsilon,1-\epsilon]$.
Each density has a positive minimum there. An elementary Legendre
evaluation estimate for an interval $I=[a,b]$ of length $\ell$ is
\begin{equation}\label{eq:legendre}
 |P(t_*)|^2\leq\frac{n^2}{\ell}R^{2(n-1)}\int_I P(t)^2\dd t
       \quad(\deg P<n),
\end{equation}
where $R=1$ if $t_*\in I$, and otherwise
$R=|z_*|+\sqrt{z_*^2-1}$, $z_*=(2t_*-a-b)/\ell$.
To see this, expand in the orthonormal basis
$\sqrt{(2j+1)/\ell}\,L_j((2t-a-b)/\ell)$ and apply
Cauchy--Schwarz. The bounds $|L_j(z)|\leq1$ on $[-1,1]$ and
$|L_j(z)|\leq(|z|+\sqrt{z^2-1})^j$ outside follow from
\[
 L_j(z)=\frac1\pi\int_0^\pi
                (z+\sqrt{z^2-1}\cos v)^j\dd v,
\]
using the imaginary square root inside the interval and parity for
negative $z$. The formula is a polynomial identity by expansion.
At $t_*=-1$, the constraint $P(-1)=1$ turns~\eqref{eq:legendre}
into a lower bound on the integral. Here
$z_*=-3/(1-2\epsilon)$. First letting $n\to\infty$
and then $\epsilon\downarrow0$ gives the matching lower exponent
for both $\lambda_{n,s}$.

Finally use $I=[0,1-\epsilon]$, where $\eta_2\geq\epsilon/8$.
The same evaluation estimate, uniformly over $t_*\in[0,1]$, gives
\[
 \|P_{n,0}\|_{[0,1]}^2\leq
 \frac{8n^2S_\epsilon^{2(n-1)}}{\epsilon(1-\epsilon)}\lambda_{n,2},
 \quad
 S_\epsilon=s_\epsilon+\sqrt{s_\epsilon^2-1},\quad
 s_\epsilon=\frac{1+\epsilon}{1-\epsilon}.
\]
Integrating against $\eta_4$ bounds $E_n$ by $\theta$ times this
quantity. Now $S_\epsilon\to1$ as $\epsilon\downarrow0$; combined
with $E_n\geq\lambda_{n,4}$ this proves the asserted rate.
\end{proof}

Write $r_n'(0)=p_n/A_n$ in lowest terms, $A_n>0$. Its actual
primitive form is
\begin{equation}\label{eq:variation-budget}
 L_n=A_n\theta-p_n=A_nE_n,
 \qquad \log L_n=\log A_n-2n\log\rho+o(n).
\end{equation}
Along any proposed subsequence, $\limsup n^{-1}\log A_n<2\log\rho$
is sufficient for shrinking values. The weak inequality is necessary,
but equality requires control of the lower-order terms. No useful
all-index reduced-denominator estimate has been established. The
quadratic logarithmic decay of the raw determinant cannot be used as
the scalar error rate: its target coefficient carries the same scale.

At $n=2$, the first-variation pair is $50\theta-49$ and the direct
weight-four pair is $82\theta-81$. The archived calculations cover
exactly the direct and first-variation forms at $n=1,\ldots,12$,
with full gcd removal. All twenty forms at $n=3,\ldots,12$ exceed one;
the four at $n=1,2$ are smaller than one. Some diagnostic logarithms are
shown below, rounded from the archived certified evaluations.
\begin{center}
\begin{tabular}{rrrr}
\toprule
$n$&$\log A_n$&$\log E_n$&$\log L_n/n^2$\\\midrule
3&21.369976&$-8.523832$&1.427349\\
6&141.949801&$-19.363907$&3.405164\\
9&354.356711&$-30.034542$&4.003977\\
12&660.303908&$-40.661541$&4.303072\\\bottomrule
\end{tabular}
\end{center}
These are finite failures, not a growth theorem for $A_n$. An earlier
table had cleared denominators without removing the common numerator
factor; at $n=12$ that overstated $\log L_n/n^2$ as $6.271803$.
The corrected data above are the retained conclusion. The archived
arcsine realization, by contrast, was identified exactly with an earlier
family; no such coefficient-sequence identification was established
for this ordinary-moment variation.
